\section{Necessity of release/refusal under one-sided verification}
\label{sec:necessity-prevented-harm}

The PARC lifecycle is motivated by the information structure of scientific
candidate queues.  The setting provides one-sided support:
\[
  A_p=1 \Rightarrow Y_p=1 ,
\]
but it does not provide the converse.  Therefore an unverified candidate cannot
be treated as negative without adding an assumption not present in the data.

\paragraph{Least-favourable null-superset principle.}
Under one-sided support, every false calibration candidate remains in the
calibration null superset.  Any method that removes unverified candidates from
the denominator must justify an additional assumption about their validity.  In
the absence of such an assumption, the conservative null-superset denominator
is the least-favourable release-card denominator.

\paragraph{Refusal lower bound.}
For requested budget \(K\), target \(\alpha\), and candidate e-values \(E_p\),
a compatible non-empty release \(R\) must satisfy
\[
  \min_{p\in R} E_p \ge {K \over \alpha |R|}.
\]
When no compatible non-empty set satisfies this inequality, refusal is the
evidence-supported output.  This is not a statement that no scientifically
valid candidates exist; it is a statement that the current release card lacks
enough one-sided evidence for the requested release.

\paragraph{Active-audit gain principle.}
Targeted verification can change the release-card state by removing verified
positives from limiting null-superset block maxima.  This explains why a small
number of high-value one-sided audits can unlock release while random audits of
the same size may leave the denominator unchanged.

\paragraph{Prevented-harm interpretation.}
False-release control is not only a statistical object.  It prevents concrete
scientific artifacts from being polluted: false lineage edges in cell tracking,
unstable crystal follow-ups in materials screening, and false persistence links
in temporal mapping.  Phase83 reports these quantities only by reformatting
completed artifacts; it does not add new empirical labels or new DFT evidence.
